problem:  
Let $\Sigma^k$ be an exotic $k$‑sphere ($k \ge 7$). For a fixed integer $r > 0$, consider a smooth rank‑$r$ vector bundle $E \to \Sigma^k$ and suppose its total space admits a complete Riemannian metric $g$ with nonnegative sectional curvature. The Cheeger–Gromoll Soul Theorem ensures that the soul $S$ of $(E,g)$ is diffeomorphic to $\Sigma^k$.  

**Problem:**  
For which pairs $(\Sigma^k, r)$ can such a metric exist? In particular, does there exist a positive integer $r$ such that some nontrivial rank‑$r$ vector bundle over an exotic sphere $\Sigma^k$ admits a complete nonnegatively curved metric whose soul $S$ is diffeomorphic to $\Sigma^k$?  
Equivalently, can exotic smooth structures on $S^k$ occur as the soul of a complete, nonnegatively curved manifold of codimension $r$, and if so, does there exist a minimal such $r$?  

Why is it a "good" problem:  
This question is mathematically consistent and precisely targets the interplay between nonnegative curvature and exotic smooth structures. It formalizes, in a sharply posed way, a geometric obstruction problem: can an exotic structure on a sphere survive as a soul under nonnegative curvature? By fixing the codimension $r$, the problem captures a progression that bridges **global Riemannian geometry** (through curvature and the soul theorem) and **differential topology** (through exotic spheres and vector bundle classifications).  
Its difficulty lies in the subtle curvature constructions on vector bundles (studied by Grove–Ziller, Belegradek–Kapovitch, and others) and the sensitivity of exotic smooth structures to bundle dimensions. Resolving even a single instance—positive or negative—would shed profound light on how curvature constrains smooth structures and could redefine our understanding of the geometric detectability of exotic differentiable types. The problem is *simple to state, deeply open, and sits precisely at the confluence of curvature, topology, and smooth classification theory*.